\documentclass[10pt,a4paper]{amsart}
\usepackage[margin=19mm]{geometry}
\usepackage{amsmath,amssymb,mathtools}
\usepackage[scr=rsfs,cal=euler]{mathalfa}
\usepackage{xfrac}
\usepackage[unicode,hidelinks,bookmarks=false]{hyperref}
\newcommand{\ie}{i.e.\ }
\newcommand{\cf}{cf.\ }
\newcommand{\resp}{respectively\ }
\newcommand{\Subsection}{\subsection}
\providecommand{\nobreakdash}{\nobreak\hbox{-}}
\setlength{\emergencystretch}{2em}
\multlinegap0pt
\usepackage{xcolor}
\newtheorem{thm}{Theorem}
\title{The Rhaly Operators on K\"othe Spaces}
\author{Nazl{\i} Do\u{g}an}
\date{Source: arXiv:2507.20214v3 (CC BY 4.0)}
\begin{document}
\maketitle

\noindent\emph{Source and licence.} The Rhaly Operators on K\"othe Spaces. Version arXiv:2507.20214v3 (CC BY 4.0), distributed by arXiv under the Creative Commons Attribution 4.0 International licence, http://creativecommons.org/licenses/by/4.0/, as recorded in the arXiv licence metadata for this exact version and shown on the abstract page https://arxiv.org/abs/2507.20214v3. Full licence notice: \texttt{licenses/arxiv-2507.20214-CC-BY-4.0.txt}.\par\smallskip

\noindent\textit{Editorial scope.} Counted unit: the article's thm TII (source0.tex lines 538--553). The article's complete original proof of it is delivered with it (source0.tex lines 554--587, one \texttt{proof} environment of 34 lines). No mathematics is written, repaired or re-derived: the statement and the proof are the article's own lines, in the article's own notation, and no retained line is substituted or re-flowed.  No further source-local context is needed: every cross-reference of every retained line resolves inside the two delivered spans.  Nothing was omitted from any retained span, so the package carries no omission record.  No retained line contains a citation, so the wrapper carries no bibliography and no bibliography entry.  No retained line contains a figure environment, an image-inclusion command or drawing code, so no image is delivered and the package declares no asset dependency.  Editorial formatting only, in addition: an AMS \texttt{amsart} preamble that restates the article's own declarations and macros used by the retained lines, namely the environments \texttt{thm} on separate counters, together with the standard packages those lines need.  The retained environments print in this document as Theorem 1 in order of appearance, whereas the article prints the same items with the numbering of its own full document.  The complete native source of the article as distributed in the arXiv e-print for this exact version is bundled unchanged as \texttt{sources/182270/source0.tex}.
\begin{thm}\label{TII}
Let $G$ denote either $\mathbb{C}$ or $\mathbb{D}$. Let $R: H(G)\to H(G)$ be a continuous linear operator. 
Then the following are equivalent:
\begin{itemize}
\item[(i)] The matrix of $R$ with respect to the Schauder basis 
$\{z^{n}\}_{n\in\mathbb{N}_{0}}$ is a Rhaly matrix, that is, there exists 
a sequence $\theta=(\theta_{n})_{n\in\mathbb{N}_{0}}$ such that
$$\hspace{1.6in}R(z^{m})=\sum^{\infty}_{j=m}\theta_{j-1}z^{j} \qquad \hspace{1.2in}m\in \mathbb{N}_{0}.$$
\item[(ii)] There exists a function $g\in H(G)$ such that $R=R_{g}$.
\end{itemize}
Moreover, in this case the sequence $\theta$ and the function $g$ are 
related by
$$\hspace{1.6in}\theta_{n}=\frac{1}{2\pi i} \int_{|w|=r} \frac{g(w)}{w^{n+1}}dw 
\qquad\hspace{1in} n\in \mathbb{N}_{0}$$
where r is a constant with $0<r<1$.
\end{thm}

\begin{proof} Let $\displaystyle f(z)=\sum^{\infty}_{n=0} a_{n}z^{n}$ and $\displaystyle g(z)=\sum^{\infty}_{n=0}b_{n}z^{n}$ be the elements of $H(G)$ We show that
$$(R_{g}f)(z)=\frac{1}{2\pi i}\int_{|w|=r_{0}} \frac{f(w)}{w(1-w)}g\left(\frac{z}{w}\right)dw=\sum^{\infty}_{n=0}\left(\sum^{n}_{k=0}a_{k}\right)b_{n}z^{n}$$
for every $z\in \mathbb{C}$. 

Let us define
$\displaystyle h(z)=\sum^{\infty}_{n=0}\left(\sum^{n}_{k=0}a_{k}\right)z^{n}$ for every $z\in \mathbb{D}$ and say $\displaystyle c_{n}=\sum^{n}_{k=0}a_{k}$ for every $n\in \mathbb{N}$. Then we have $\displaystyle h(z)=\sum^{\infty}_{n=0}c_n z^n$. By the following of equalities
\begin{equation*}
\begin{split}
h(z)&= \sum^{\infty}_{n=0} \left(\sum^{n}_{k=0}a_{k}\right)z^{n}=\sum^{\infty}_{k=0}a_{k}\left(\sum^{\infty}_{j=k} z^{j}\right) \\
&=\sum^{\infty}_{k=0} a_{k}z^{k}\left(\sum^{\infty}_{j=0}z^{j}\right)=\frac{1}{1-z}\sum^{\infty}_{k=0}a_{k}z^{k}=\frac{1}{1-z}f(z),
\end{split}
\end{equation*}
hence $\displaystyle h(z)=\frac{f(z)}{1-z}$ for every $z\in \mathbb{D}$. 

Fix a $z\in G$ and a 
$r_{0}$ satisfying $|z|<r_{0}$ and $r_{0}\neq 1$.  We have
\begin{equation*}
\begin{split}
&\frac{1}{2\pi i}\int_{|w|=r_{0}} \frac{f(w)}{w(1-w)}g\left(\frac{z}{w}\right)dw=\frac{1}{2\pi i}\int_{|w|=r_{0}} \frac{h(w)}{w}g\left(\frac{z}{w}\right)dw \\
&= \sum^{\infty}_{n=0}b_{n}z^{n} \frac{1}{2\pi i}\int_{|w|=r_{0}} \frac{h(w)}{w^{n+1}}dw=\sum^{\infty}_{n=0}c_nb_nz^n=\sum^{\infty}_{n=0}\left(\sum^{n}_{k=0}a_{k}\right)b_{n}z^n
\end{split}
\end{equation*}
for every $z\in \mathbb{C}$. Consequently,
$$(R_{g}z^{n})(\xi)=\sum^{\infty}_{m=n}b_{m}\xi^{m}$$
for every $n\in \mathbb{N}_{0}$. Hence, the associated matrix of $R_{g}$ is a Rhaly matrix and 
$$\theta_{n}=b_{n}=\frac{1}{2\pi i} \int_{|w|=1} \frac{g(w)}{w^{n+1}}dw$$
for every $n\in \mathbb{N}_{0}$.

On the other hand, if $R:H(G)\to H(G)$ is a continuous linear operator whose associated matrix is a Rhaly matrix, the function $\displaystyle g(z)=\sum^{\infty}_{n=0}\theta_n z^n$, formed by the coefficients of the sequence $\theta$, is an element of $H(G)$ since
\begin{equation*}
(R1)(z)=\sum^{\infty}_{n=0}\theta_{n}z^{n}=g(z)\in H(G). 
\end{equation*}
The above calculation likewise makes it evident that $R=R_{g}$.
\end{proof}

\end{document}
